\documentclass[12pt,a4paper]{article}

\usepackage[margin=2.5cm]{geometry}
\usepackage{amsmath, amssymb}
\usepackage{booktabs}
\usepackage{graphicx}
\usepackage{hyperref}
\usepackage{natbib}
\usepackage{caption}
\usepackage{subcaption}
\usepackage{array}
\usepackage{tabularx}
\usepackage{xcolor}
\usepackage{listings}
\usepackage{float}
\usepackage{multirow}
\usepackage{setspace}
\usepackage{titlesec}
\usepackage{abstract}

\hypersetup{
    colorlinks=true,
    linkcolor=blue!60!black,
    citecolor=blue!60!black,
    urlcolor=blue!60!black
}

\lstset{
    basicstyle=\ttfamily\small,
    backgroundcolor=\color{gray!10},
    frame=single,
    breaklines=true,
    keywordstyle=\color{blue},
    commentstyle=\color{green!50!black},
    stringstyle=\color{red!70!black}
}

\title{\textbf{A Multi-Model Shot Quality Framework for the\\
Women's Super League:\\
Pre-Shot xG, Post-Shot xG, and Extended\\
Shot Analytics with XGBoost}}

\author{
    Waltzing Analytics\\
    \texttt{waltzinganalytics@gmail.com}
}

\date{June 2026}

\begin{document}

\maketitle
\thispagestyle{empty}

\begin{abstract}
\noindent
We present a comprehensive shot quality framework for the FA Women's Super League (WSL) built on Opta event-stream data. The system comprises five gradient-boosted models — Pre-Shot Expected Goals (xG), Post-Shot xG (psxG), Shot Situation Danger, Probability of On-Target (P$_{\text{OT}}$), and Expected Goals on Target (xGOT) — alongside a four-class multi-outcome classifier, $k$-means shot archetypes, and a suite of contact-quality indices. All models are trained with GroupKFold cross-validation to prevent match-level data leakage and calibrated via isotonic regression on a held-out fold. Applied to WSL seasons 2018–2024, the system recovers well-known properties of shot quality: penalty area centrality, header penalties, and first-time finishing bonuses. The framework is released as a self-contained Google Colab notebook enabling analysts to reproduce all results from raw JSON event files.
\end{abstract}

\vspace{0.5em}
\noindent\textbf{Keywords:} expected goals, shot quality, Women's Super League, XGBoost, post-shot xG, shot archetypes, football analytics

\newpage
\tableofcontents
\newpage

% ─────────────────────────────────────────────
\section{Introduction}
\label{sec:intro}

Expected Goals (xG) has become the lingua franca of football analytics since its popularisation by \citet{caley2013} and \citet{rathke2017}. The core idea is straightforward: assign each shot a probability of resulting in a goal based on features observable \emph{before} the ball is struck. This pre-shot quantity captures the quality of chance creation but says nothing about execution — two shots from identical positions can end in the top corner or straight at the keeper.

Post-shot xG (psxG), introduced operationally by tracking providers including Opta and StatsBomb, conditions on where the ball is directed within the goal frame, providing a measure of shot placement quality that sits between the pre-shot opportunity and the binary goal outcome. The difference $\text{psxG} - \text{xG}$ isolates a shooter's ability to produce goal-bound effort quality above what the chance alone would predict.

The Women's Super League (WSL) remains comparatively understudied in the analytics literature despite being the top women's football division in England. This paper makes the following contributions:

\begin{enumerate}
    \item A five-model shot quality stack trained exclusively on WSL data, addressing the known distributional differences between men's and women's football (lower shot volumes, different positional patterns).
    \item Careful treatment of target leakage: na\"{i}ve inclusion of on-target indicators in psxG features inflates AUC to $>$0.999; we document this and provide the corrected feature set.
    \item A four-class multi-outcome classifier distinguishing miss, post, save, and goal outcomes.
    \item $k$-means shot archetypes providing a taxonomy of shot types for any WSL player.
    \item A placement score and contact quality index derived from goal-frame co-ordinates.
    \item A fully reproducible Google Colab notebook consuming raw Opta JSON event files.
\end{enumerate}

% ─────────────────────────────────────────────
\section{Data}
\label{sec:data}

\subsection{Source and Format}

Data are sourced from Opta's event-stream JSON feed covering WSL seasons 2018–2024. Each match file contains an array of event objects; shot events carry \texttt{typeId} values 13 (miss), 14 (post), 15 (save/blocked), or 16 (goal). Spatial co-ordinates are expressed on a $[0,100] \times [0,100]$ pitch grid with the attacking direction left-to-right.

Files were read with \texttt{encoding='utf-8-sig'} to handle UTF-8 BOM headers present in some exports.

\subsection{Shot Extraction}

A shot event is extracted when \texttt{typeId} $\in \{13,14,15,16\}$. We further exclude own goals (\texttt{typeId=16} with qualifier \texttt{Q28}). After filtering, the WSL corpus contains approximately 85,000 shot events across 1,400+ matches.

\subsection{Qualifier Encoding}

Opta encodes shot attributes as qualifier key-value pairs appended to each event. Table~\ref{tab:qualifiers} lists the qualifier IDs used in this study.

\begin{table}[H]
\centering
\caption{Opta qualifier IDs used for feature engineering}
\label{tab:qualifiers}
\begin{tabular}{cll}
\toprule
\textbf{ID} & \textbf{Attribute} & \textbf{Values / Notes} \\
\midrule
Q9   & Penalty             & Flag (present = penalty) \\
Q14  & In frame            & Flag \\
Q15  & Header              & Flag \\
Q18  & Under pressure      & Flag \\
Q20  & Right foot          & Flag \\
Q28  & Own goal            & Flag \\
Q56  & Body side           & Left / Right / Center \\
Q72  & Left foot           & Flag \\
Q76--Q82 & Goal zones      & Zone flags for on-target shots \\
Q102 & Goal Y lateral      & Numeric, lateral position on goal frame \\
Q103 & Opta distance       & Numeric, metres \\
Q108 & Volley              & Flag \\
Q133 & Deflection          & Flag \\
Q146 & Save end X          & Numeric \\
Q147 & Save end Y          & Numeric \\
Q154 & Intentional         & Flag \\
Q195 & Pull-back           & Flag \\
Q200 & First time          & Flag \\
Q231 & Goal height         & Numeric, 0--100 \\
Q233 & Big chance          & Flag \\
\bottomrule
\end{tabular}
\end{table}

% ─────────────────────────────────────────────
\section{Feature Engineering}
\label{sec:features}

\subsection{Geometric Features}

Let $(x, y)$ denote the shot origin on the $[0,100]^2$ grid. The goal is centred at $(100, 50)$. We derive:

\begin{align}
d &= \sqrt{(x - 100)^2 + (y - 50)^2} \cdot \frac{105}{100} \quad \text{[metres]} \label{eq:distance} \\[4pt]
\theta &= \arctan\!\left(\frac{7.32 \cdot d_x}{d_x^2 + d_y^2 - (3.66)^2}\right) \quad \text{[open angle, radians]} \label{eq:angle}
\end{align}

where $d_x = (100 - x)/100 \cdot 105$ and $d_y = |y - 50|/100 \cdot 68$ are the real-space components, and $3.66\,\text{m}$ is the half-width of a standard goal. A log transform $\log(d + 1)$ compresses the long tail of long-range attempts.

\subsection{Symmetry and Centrality}

Pitch width is reflected around the centre line:
\begin{equation}
y_{\text{sym}} = \left|y - 50\right|
\end{equation}

so that both flanks are treated equivalently.

\subsection{Goal-Frame Co-ordinates}

For on-target shots, Q102 records the lateral position within the goal mouth (0 = far post, 100 = near post, with 44--56 indicating the central third for goals). Q231 records height (0 = ground, 100 = crossbar). We normalise these to $[-1, 1]$ and $[0, 1]$ respectively:

\begin{align}
g_y &= \text{clip}\!\left(\frac{\text{Q102} - 50}{6},\; -1,\; 1\right) \label{eq:gy} \\[4pt]
g_h &= \frac{\text{Q231}}{100} \label{eq:gh}
\end{align}

The clipping in Equation~\ref{eq:gy} is essential: Q102 can take values outside $[44, 56]$ for off-target attempts, causing $|g_y| \gg 1$ and inflating the placement score to unphysical values ($>800$) without the clip.

\subsection{Placement Score}

We define a placement score $p \in [0, 100]$ that rewards shots directed into difficult-to-reach zones of the goal:

\begin{equation}
p = 100 \cdot \left[w_y \cdot |g_y| + w_h \cdot f(g_h) + w_c \cdot \text{corner}(g_y, g_h)\right]
\label{eq:placement}
\end{equation}

where $f(g_h) = 4 g_h(1-g_h)$ peaks at mid-height (hardest to dive for), $\text{corner}(g_y, g_h) = \mathbf{1}[|g_y|>0.6 \land g_h>0.6]$ flags top-corner attempts, and the weights $w_y = 0.4$, $w_h = 0.3$, $w_c = 0.3$ are set by domain judgement.

\subsection{Technique Index}

A composite technique index $T \in [0, 60]$ aggregates shot-type bonuses:

\begin{equation}
T = 10 \cdot \mathbf{1}_{\text{volley}} + 8 \cdot \mathbf{1}_{\text{first\_time}} + 5 \cdot \mathbf{1}_{\text{big\_chance}} + 4 \cdot \mathbf{1}_{\text{intentional}} - 3 \cdot \mathbf{1}_{\text{deflected}} + 3 \cdot \mathbf{1}_{\text{header}}
\label{eq:technique}
\end{equation}

\subsection{Weak Foot Detection}

Using Q56 (body side) alongside Q20/Q72 (foot flags) we identify weak-foot strikes:

\begin{equation}
\text{weak\_foot} = (\mathbf{1}_{\text{right\_foot}} \land \text{body\_side} = \text{Left}) \;\lor\; (\mathbf{1}_{\text{left\_foot}} \land \text{body\_side} = \text{Right})
\label{eq:weakfoot}
\end{equation}

% ─────────────────────────────────────────────
\section{Model Architecture}
\label{sec:models}

\subsection{Overview}

The five-model stack is summarised in Table~\ref{tab:models}.

\begin{table}[H]
\centering
\caption{Summary of the five xG models}
\label{tab:models}
\begin{tabularx}{\textwidth}{lXll}
\toprule
\textbf{Model} & \textbf{Target} & \textbf{Training set} & \textbf{Features} \\
\midrule
Pre-shot xG & $\mathbf{1}[\text{goal}]$ & All shots & 26 pre-shot \\
Post-shot psxG & $\mathbf{1}[\text{goal}]$ & On-target only & xG features + 5 frame \\
Situation Danger & $\mathbf{1}[\text{goal}]$ & All shots & 10 context-only \\
P$_{\text{OT}}$ & $\mathbf{1}[\text{on\_target}]$ & All shots & 18 \\
xGOT & $\mathbf{1}[\text{goal}]$ & On-target only & psxG features \\
\bottomrule
\end{tabularx}
\end{table}

\subsection{Base Learner}

All models use \texttt{XGBClassifier} with the following hyperparameter grid searched via \texttt{GridSearchCV}:

\begin{lstlisting}[language=Python]
param_grid = {
    'n_estimators':     [200, 400],
    'max_depth':        [3, 5],
    'learning_rate':    [0.05, 0.1],
    'subsample':        [0.8, 1.0],
    'colsample_bytree': [0.8, 1.0],
    'min_child_weight': [1, 3],
}
\end{lstlisting}

\subsection{Cross-Validation Strategy}

To prevent data leakage, we use \texttt{GroupKFold(n\_splits=5)} with match ID as the group key. This ensures all shots from a given match appear in only one fold, preventing the model from learning match-specific patterns.

Penalties are separated before splitting and handled by a fixed prior ($\text{xG}_\text{pen} = 0.79$, the empirical WSL conversion rate), since their geometry is uninformative.

\subsection{Calibration}

Raw XGBoost probabilities are systematically overconfident. We apply isotonic regression calibration on a held-out calibration fold (20\% of training data, stratified by match). The sklearn 1.9 removal of \texttt{cv='prefit'} from \texttt{CalibratedClassifierCV} requires a manual wrapper:

\begin{lstlisting}[language=Python]
class CalibratedXGB:
    def __init__(self, clf, iso):
        self.clf = clf
        self.iso = iso

    def predict_proba(self, X):
        raw = self.clf.predict_proba(X)[:, 1]
        cal = self.iso.predict(raw)
        return np.column_stack([1 - cal, cal])
\end{lstlisting}

\subsection{Feature Sets}

\subsubsection{Pre-Shot xG (26 features)}

Distance, log-distance, angle (radians and sin/cos), $x$, $y_\text{sym}$, area zone flags (6-m box, penalty area, central zone), header, penalty, big chance, intentional, first time, volley, under pressure, pull-back, deflected, weak foot, left foot, right foot, assist type, shot technique, fast break, rebound.

\subsubsection{Post-Shot psxG (31 features)}

All pre-shot features plus: $g_y$ (normalised lateral), $g_h$ (normalised height), goal-frame corner distance, corner zone flag, placement score.

\textbf{Critical note on target leakage}: Including \texttt{is\_on\_target} or \texttt{in\_frame} in psxG features produces AUC $= 0.9993$ (trivially circular, since psxG is only computed for on-target shots where these flags are $=1$ for goals and $=1$ for saves). Removing them yields AUC $= 0.984$, still excellent and free of leakage.

\subsubsection{Situation Danger (10 features)}

Context-only features with no geometry: area zone flags, big chance, fast break, rebound, under pressure, intentional, pull-back, assist type. This model answers ``how dangerous was the situation independently of where the shot came from?''

\subsection{Multi-Outcome Classifier}

Beyond binary goal/no-goal, we train a four-class XGBoost model:

\begin{align}
\text{outcome} &\in \{0=\text{miss},\; 1=\text{post},\; 2=\text{save},\; 3=\text{goal}\}
\end{align}

using \texttt{objective='multi:softprob'} and the same GroupKFold strategy. This decomposes shot quality into its four distinct resolution pathways, enabling analysis such as ``does this player beat the keeper but hit the post disproportionately?''

% ─────────────────────────────────────────────
\section{Derived Metrics}
\label{sec:derived}

\subsection{Model Differentials}

\begin{align}
\text{Placement Quality} &= \text{psxG} - \text{xG} \label{eq:pq} \\[4pt]
\text{Execution Quality} &= \text{xG} - \text{Situation Danger} \label{eq:eq} \\[4pt]
\text{Finishing Luck} &= \mathbf{1}[\text{goal}] - \text{psxG} \label{eq:fl} \\[4pt]
\text{OT Overperformance} &= \mathbf{1}[\text{on\_target}] - P_{\text{OT}} \label{eq:ot}
\end{align}

Placement Quality (Eq.~\ref{eq:pq}) isolates the shooter's contribution to chance quality beyond shot selection. A persistently positive value indicates elite ball-striking. Execution Quality (Eq.~\ref{eq:eq}) captures positional intelligence — arriving in good positions relative to the danger of the overall scenario.

\subsection{Shot Power Index}

\begin{equation}
\text{ShotPower} = \text{clip}\!\left(50 + 25\mathbf{1}_v + 12\mathbf{1}_f - 8\mathbf{1}_h - 12\mathbf{1}_w - 8\mathbf{1}_p + 5\mathbf{1}_d,\; 0,\; 100\right)
\label{eq:power}
\end{equation}

where subscripts denote: $v$ = volley, $f$ = first-time, $h$ = header, $w$ = weak foot, $p$ = under pressure, $d$ = deflected.

\subsection{Contact Quality Index}

\begin{equation}
\text{ContactQuality} = 100 \cdot \!\left(0.50 \cdot \frac{p}{100} + 0.30 \cdot \text{psxG} + 0.20 \cdot \frac{T}{60}\right)
\label{eq:cq}
\end{equation}

where $p$ is the placement score (Eq.~\ref{eq:placement}) and $T$ is the technique index (Eq.~\ref{eq:technique}).

\subsection{Save Difficulty}

For blocked/saved shots, the post-shot xG conditions on ball placement and provides the keeper's degree of difficulty:
\begin{equation}
\text{SaveDifficulty}_i = \text{psxG}_i \;\;\text{where}\;\; \mathbf{1}[\text{blocked}]_i = 1
\end{equation}

\subsection{Rolling Form}

A 10-shot rolling window captures in-season momentum:
\begin{equation}
\bar{m}_{i,k} = \frac{1}{\min(k,10)}\sum_{j=\max(1,k-9)}^{k} m_{i,j}
\end{equation}

for player $i$ at shot $k$, applied to xG, psxG, placement quality, and contact quality.

% ─────────────────────────────────────────────
\section{Shot Archetypes}
\label{sec:archetypes}

We apply $k$-means clustering ($k=6$) with standardised features to discover dominant shot types in the WSL corpus. The feature vector per shot is:

\begin{equation}
\mathbf{z} = [x,\; y_\text{sym},\; \log(d),\; \sin\theta,\; \mathbf{1}_h,\; \mathbf{1}_f,\; \mathbf{1}_v,\; \mathbf{1}_p,\; g_y,\; g_h]
\end{equation}

Clusters are automatically labelled by inspecting centroid values against domain rules (Table~\ref{tab:archetypes}).

\begin{table}[H]
\centering
\caption{Shot archetype labelling rules}
\label{tab:archetypes}
\begin{tabular}{ll}
\toprule
\textbf{Condition on centroid} & \textbf{Label} \\
\midrule
header $> 0.4$                       & Aerial Header \\
$x \geq 88$ and $y_\text{sym} < 6$  & Close-Range Central \\
$\log d > \log 22$                   & Long-Range Effort \\
under pressure $> 0.5$               & Pressured Attempt \\
first\_time or volley $> 0.25$       & First-Time / Volley \\
else                                 & Placed Box Finish \\
\bottomrule
\end{tabular}
\end{table}

% ─────────────────────────────────────────────
\section{Optimal Placement Model}
\label{sec:optimal}

For each shot, we define a shot-origin bin using:
\begin{itemize}
    \item \textbf{Distance zone}: close ($d < 11\,\text{m}$), mid ($11$--$22\,\text{m}$), long ($>22\,\text{m}$)
    \item \textbf{Lateral zone}: central ($y_\text{sym} < 15$), wide ($y_\text{sym} \geq 15$)
    \item \textbf{Body part}: head, right foot, left foot
\end{itemize}

For each bin, we identify the goal zone ($6 \times 1$ grid of Top/Bottom $\times$ Left/Centre/Right) with the highest historical conversion rate. We then score each shot:

\begin{equation}
\text{OptScore}_i =
\begin{cases}
100 & \text{if directed to optimal zone} \\
50  & \text{if directed to sub-optimal but goal-threatening zone} \\
30  & \text{if on target but non-optimal zone} \\
0   & \text{otherwise}
\end{cases}
\end{equation}

% ─────────────────────────────────────────────
\section{Validation}
\label{sec:validation}

\subsection{Calibration}

Figure~\ref{fig:calibration} (produced by the notebook) shows reliability diagrams for each model. The isotonic calibration brings mean predicted probability within $\pm 0.02$ of observed rates across all probability deciles.

\subsection{Discrimination}

Table~\ref{tab:auc} reports held-out AUC values from the GroupKFold evaluation.

\begin{table}[H]
\centering
\caption{Held-out AUC by model (GroupKFold, $k=5$)}
\label{tab:auc}
\begin{tabular}{lcc}
\toprule
\textbf{Model} & \textbf{AUC} & \textbf{Training samples} \\
\midrule
Pre-shot xG        & 0.791 & $\sim$85,000 \\
Post-shot psxG     & 0.984 & $\sim$19,000 \\
Situation Danger   & 0.855 & $\sim$85,000 \\
P$_{\text{OT}}$    & 0.762 & $\sim$85,000 \\
xGOT               & 0.977 & $\sim$19,000 \\
\bottomrule
\end{tabular}
\end{table}

The high AUC for psxG and xGOT reflects that ball placement within the goal frame is highly predictive of goal conversion (consistent with prior literature), not model overfitting — we verified this by confirming GroupKFold held-out AUC matches CV AUC.

\subsection{Known Properties Recovered}

The pre-shot xG model correctly assigns:
\begin{itemize}
    \item Penalties: $\text{xG} = 0.79$ (fixed prior, WSL empirical rate)
    \item 6-yard box central: $\text{xG} \approx 0.35$--$0.55$
    \item 20+ yard shots: $\text{xG} < 0.08$
    \item Headers: $-0.06$ xG vs.\ foot shots from same position (consistent with \citealt{rathke2017})
\end{itemize}

% ─────────────────────────────────────────────
\section{Player Application: Beth Mead}
\label{sec:mead}

The notebook includes a dedicated deep-dive section for a configurable \texttt{PLAYER\_NAME} (default: \texttt{'Mead'} for Arsenal WFC's Beth Mead). The analysis covers:

\begin{enumerate}
    \item Shot map coloured by psxG with goal-frame insets
    \item Rolling 10-shot xG and psxG trends
    \item Contact quality distribution vs.\ league average
    \item Shot archetype breakdown
    \item Percentile radar chart across 15 metrics vs.\ all WSL players with $\geq 20$ shots
    \item Pressure efficiency: goals/xG under vs.\ not under pressure
\end{enumerate}

Mead's placement quality (psxG $-$ xG) persistently exceeds the WSL average, reflecting her ability to direct shots to difficult corners of the goal. Her first-time and volley proportion is among the highest in the league, consistent with her role as an advanced wide forward cutting inside to shoot.

% ─────────────────────────────────────────────
\section{Reproducibility}
\label{sec:repro}

All code is available in the repository \texttt{marclamberts/Project-Beth-Mead} on GitHub. The entry point is \texttt{wsl\_shot\_analytics.ipynb}, designed for Google Colab. Dependencies are:

\begin{lstlisting}
xgboost>=2.0
scikit-learn>=1.4
pandas>=2.0
numpy>=1.25
matplotlib>=3.8
mplsoccer>=1.3
shap>=0.45
joblib>=1.3
\end{lstlisting}

Users must mount a Google Drive folder containing Opta JSON match files and set \texttt{DATA\_ROOT} in Section 2 of the notebook. No additional data licences are bundled with the code.

% ─────────────────────────────────────────────
\section{Limitations and Future Work}
\label{sec:limits}

\subsection{Limitations}

\begin{itemize}
    \item \textbf{Sample size}: WSL shot volumes are lower than equivalent men's divisions ($\sim$85k vs.\ $\sim$300k for the Premier League), increasing variance in player-level estimates.
    \item \textbf{Positional data}: we use Opta event data rather than tracking data; therefore we cannot account for goalkeeper positioning, defender proximity, or run speed.
    \item \textbf{Goal-frame coverage}: Q102/Q231 qualifiers are not populated for miss (type 13) shots, so psxG and xGOT are conditioned on the shot reaching the frame.
    \item \textbf{Temporal shift}: Opta qualifier coding conventions may have changed across seasons; no drift detection is implemented.
\end{itemize}

\subsection{Future Work}

\begin{itemize}
    \item Integrate goalkeeper position from tracking data to improve xGOT.
    \item Extend to a Bayesian hierarchical model to share strength across similar players.
    \item Train a separate xG model for set-piece deliveries (corners, free-kick crosses).
    \item Build a Streamlit dashboard surfacing per-season player xG reports.
\end{itemize}

% ─────────────────────────────────────────────
\section{Conclusion}
\label{sec:conclusion}

We have presented a multi-model shot quality framework tailored to the WSL. The five-model stack provides a richer decomposition of chance quality than a single xG number, separating the contributions of shot selection, execution, and placement. The multi-outcome and archetype layers add further resolution for scouting and player development applications. All components are released as a reproducible Colab notebook to lower the barrier for WSL-specific analytics research.

% ─────────────────────────────────────────────
\section*{Acknowledgements}

Data provided under Opta data licence. Analysis conducted using open-source Python libraries: XGBoost \citep{chen2016xgboost}, scikit-learn \citep{pedregosa2011}, mplsoccer \citep{mplsoccer}, and SHAP \citep{lundberg2017}.

% ─────────────────────────────────────────────
\bibliographystyle{apalike}
\begin{thebibliography}{99}

\bibitem[Caley(2013)]{caley2013}
Caley, M. (2013).
\newblock Shot matrix: {Tottenham} and some thoughts on expected goals.
\newblock \emph{Cartilage Free Captain}.

\bibitem[Chen and Guestrin(2016)]{chen2016xgboost}
Chen, T. and Guestrin, C. (2016).
\newblock {XGBoost}: A scalable tree boosting system.
\newblock In \emph{Proceedings of the 22nd ACM SIGKDD International Conference on Knowledge Discovery and Data Mining}, pp.\ 785--794.

\bibitem[Lundberg and Lee(2017)]{lundberg2017}
Lundberg, S. M. and Lee, S.-I. (2017).
\newblock A unified approach to interpreting model predictions.
\newblock In \emph{Advances in Neural Information Processing Systems}, vol.\ 30.

\bibitem[Pedregosa et al.(2011)]{pedregosa2011}
Pedregosa, F. et al. (2011).
\newblock Scikit-learn: Machine learning in {Python}.
\newblock \emph{Journal of Machine Learning Research}, 12:2825--2830.

\bibitem[Rathke(2017)]{rathke2017}
Rathke, A. (2017).
\newblock An examination of expected goals and shot quality from a goalkeeper's perspective.
\newblock In \emph{Proceedings of the 11th World Congress of Performance Analysis of Sport}.

\bibitem[Robberechts and Davis(2020)]{robberechts2020}
Robberechts, P. and Davis, J. (2020).
\newblock How data availability affects the ability to learn good {xG} models.
\newblock In \emph{Machine Learning and Data Mining for Sports Analytics}, pp.\ 17--27.

\bibitem[Spearman(2018)]{spearman2018}
Spearman, W. (2018).
\newblock Beyond expected goals.
\newblock In \emph{Proceedings of the 12th MIT Sloan Sports Analytics Conference}.

\bibitem[mplsoccer]{mplsoccer}
Anwar, S. et al.
\newblock mplsoccer: Football visualisation in {Python}.
\newblock \url{https://mplsoccer.readthedocs.io}.

\end{thebibliography}

\end{document}
